% preamble-exam-zh.tex — 极简讲解页共享 preamble
% 目标：复刻「题目独立、提示轻框、路线箭头、主体双栏、答案轻收束」的讲义风格。

\documentclass{exam-zh}

% ---------- 基础配置 ----------
\examsetup{
  page/size = a4paper,
  page/show-head = false,
  page/show-foot = false,
  sealline/show = false,
  font = none,
  math-font = none,
}

% ---------- 包 ----------
\usepackage{geometry}
\geometry{top=10mm,bottom=14mm,left=15mm,right=13mm}
\AtBeginDocument{\newgeometry{top=10mm,bottom=14mm,left=15mm,right=13mm}}

\usepackage{xcolor}
\usepackage{needspace}
\usepackage{enumitem}
\usepackage{array}
\usepackage{tabularx}
\usepackage{tcolorbox}
\tcbuselibrary{skins,breakable}
\usepackage{paracol}

% ---------- 打印/屏幕颜色切换 ----------
% 用法：\documentclass[monochrome]{exam-zh} 可降级为灰度。
% 注意：不要使用 \if@xxx 放在 makeatother 外部；这里改成普通条件名。
\newif\ifeduprintcolor
\eduprintcolortrue
\makeatletter
\@ifclasswith{exam-zh}{monochrome}{\eduprintcolorfalse}{}
\makeatother

% ---------- 语义颜色 ----------
\ifeduprintcolor
  \definecolor{edu-blue}{RGB}{31,62,126}
  \definecolor{edu-blue-soft}{RGB}{232,238,250}
  \definecolor{edu-blue-bg}{RGB}{248,250,254}
  \definecolor{edu-ink}{RGB}{30,35,45}
  \definecolor{edu-muted}{RGB}{118,124,136}
  \definecolor{edu-line}{RGB}{209,218,235}
  \definecolor{edu-soft-bg}{RGB}{250,251,253}
  \definecolor{edu-red}{RGB}{168,58,58}
  \definecolor{edu-red-bg}{RGB}{255,247,245}
\else
  \definecolor{edu-blue}{gray}{0.15}
  \definecolor{edu-blue-soft}{gray}{0.90}
  \definecolor{edu-blue-bg}{gray}{0.98}
  \definecolor{edu-ink}{gray}{0.10}
  \definecolor{edu-muted}{gray}{0.45}
  \definecolor{edu-line}{gray}{0.82}
  \definecolor{edu-soft-bg}{gray}{0.97}
  \definecolor{edu-red}{gray}{0.28}
  \definecolor{edu-red-bg}{gray}{0.97}
\fi

% ---------- 全局排版 ----------
\setlength{\parindent}{0pt}
\setlength{\parskip}{0pt}
\linespread{1.16}
\renewcommand{\arraystretch}{1.15}

% ctex 通常提供 \kaishu；这里用它做教师旁白感。
\newcommand{\edukai}{\kaishu}

% ---------- 顶部元信息 ----------
\newcommand{\eduTopBar}[2]{%
  \noindent
  \begin{tabularx}{\linewidth}{@{}Xr@{}}
    {\color{edu-blue}\bfseries #1} &
    {\color{edu-blue}\bfseries #2}
  \end{tabularx}
  \vspace{8mm}
}

% ---------- 轻标题体系 ----------
% 只有真正的大结构才用它；不要给提示/路线滥用标题。
\newcommand{\eduSectionTitle}[1]{%
  \needspace{5\baselineskip}%
  \vspace{2mm}
  {\color{edu-blue}\bfseries\large #1}\par
  \vspace{3mm}
}

\newcommand{\eduSubTitle}[1]{%
  \needspace{4\baselineskip}%
  \vspace{1mm}
  {\color{edu-blue}\bfseries #1}\par
  \vspace{1mm}
}

% ---------- 小问题干：复用 exam (1)(2)(3) 格式 ----------
\newenvironment{eduSubQuestionItem}[1]{%
  \needspace{4\baselineskip}%
  \vspace{2mm}
  \par\noindent
  \hangindent=8mm\hangafter=1
  {\color{edu-blue}\bfseries #1}\enspace
  \begingroup
  \color{edu-ink}
}{%
  \endgroup
  \par
  \vspace{3mm}
}

% ---------- 辅助信息条（解题流程/关键想法/思考）楷体 ----------
\newenvironment{eduAuxItem}[1]{%
  \needspace{4\baselineskip}%
  \vspace{2mm}
  \par\noindent
  \hangindent=8mm\hangafter=1
  {\color{edu-blue}\edukai $\triangleright$~#1}\enspace
  \begingroup
  \color{edu-ink}
  \edukai
}{%
  \endgroup
  \par
  \vspace{4mm}
}

\newcommand{\eduColumnTitle}[2]{%
  {\color{edu-blue}\bfseries #1\hspace{1.5mm}#2}\par
  \vspace{2.5mm}
}

% ---------- 图标：不用外部 icon 字体，保证可移植 ----------
\newcommand{\eduIconBulb}{\(\circ\)}
\newcommand{\eduIconBook}{\(\square\)}
\newcommand{\eduIconCheck}{\(\checkmark\)}
\newcommand{\eduIconPin}{\(\diamond\)}
\newcommand{\eduIconNote}{\(\triangleright\)}

% ---------- 题目区 ----------
% 题目区是页面入口：弱标签 + 一条长蓝线 + 大题干。
\newenvironment{eduProblemBlock}[1][题目]{%
  \needspace{8\baselineskip}
  {\small\color{edu-muted}#1}\par
  \vspace{1.5mm}
  {\color{edu-blue}\hrule height 0.55pt}
  \vspace{4.5mm}
  \begingroup
  \color{edu-ink}
  \large
  \setlength{\parskip}{2.2mm}
  \linespread{1.20}\selectfont
}{%
  \par
  \endgroup
  \vspace{6mm}
}

\newenvironment{eduSubQuestions}{%
  \begin{enumerate}[
    label=(\arabic*),
    leftmargin=8mm,
    itemsep=2.0mm,
    topsep=2.0mm,
    parsep=0pt
  ]
}{%
  \end{enumerate}
}

% ---------- 读题提示：轻框，不做同级标题 ----------
\newtcolorbox{eduReadTipBox}{
  enhanced,
  breakable,
  colback=white,
  colframe=edu-blue!45,
  boxrule=0.45pt,
  arc=1.6mm,
  left=10mm,
  right=4mm,
  top=5mm,
  bottom=5mm,
  before skip=1mm,
  after skip=7mm,
  fontupper=\edukai\normalsize,
  colupper=edu-ink,
  overlay={
    \node[anchor=north west, text=edu-blue] at ([xshift=3.2mm,yshift=-3.2mm]frame.north west)
    {\small\eduIconNote};
  }
}

\newenvironment{eduTipList}{%
  \begin{itemize}[
    label=\textbullet,
    leftmargin=5mm,
    itemsep=1.2mm,
    topsep=0pt,
    parsep=0pt
  ]
}{%
  \end{itemize}
}

% ---------- 解题路线：虚线框 + 文字箭头 ----------
\newenvironment{eduRouteFlow}{%
  \needspace{4\baselineskip}%
  \vspace{1mm}
  \begin{center}
  \normalsize
}{%
  \end{center}
  \vspace{7mm}
}
\newcommand{\eduRouteStep}[1]{%
  {\color{edu-ink}#1}%
}
\newcommand{\eduRouteArrow}{%
  \;$\boldsymbol{\to}$\;%
}

% ---------- 主体讲解：使用 paracol 实现跨页自适应双栏 ----------
\newenvironment{eduDualExplain}{%
  \needspace{6\baselineskip}% 防止标题孤行
  \columnratio{0.25, 0.75}%
  \setlength{\columnsep}{7mm}%
  \columnseprule=0.45pt%
  \colseprulecolor{edu-blue}%
  \begin{paracol}{2}% 开启双栏
}{%
  \end{paracol}% 结束双栏
  \vspace{5mm}
}

\newenvironment{eduExplainLeft}{%
  \normalsize\color{edu-ink}%
}{%
  \switchcolumn
}

\newcommand{\eduColumnDivider}{}

\newenvironment{eduExplainRight}{%
  \normalsize\color{edu-ink}%
}{%
}

\newenvironment{eduNumberList}{%
  \begin{enumerate}[
    label=\arabic*.,
    leftmargin=6mm,
    itemsep=4.8mm,
    topsep=0pt,
    parsep=0pt
  ]
}{%
  \end{enumerate}
}

\newenvironment{eduBulletList}{%
  \begin{itemize}[
    label=\textbullet,
    leftmargin=5mm,
    itemsep=1.2mm,
    topsep=0pt,
    parsep=0pt
  ]
}{%
  \end{itemize}
}

% ---------- 底部双栏：答案 + 方法提醒 ----------
\newenvironment{eduBottomGrid}{%
  \vspace{1mm}
  \par\noindent
}{%
  \par\vspace{1mm}
}

\newenvironment{eduSummaryLeft}{%
  \begin{minipage}[t]{0.30\linewidth}
}{%
  \end{minipage}
}

\newcommand{\eduSummaryDivider}{%
  \hspace{4mm}{\color{edu-line}\vrule width 0.35pt}\hspace{7mm}%
}

\newenvironment{eduSummaryRight}{%
  \begin{minipage}[t]{0.58\linewidth}
}{%
  \end{minipage}
}

% ---------- 答案/方法提醒：轻量收束 ----------
\newtcolorbox{eduSummaryBlock}[2]{
  enhanced,
  breakable,
  colback=white,
  colframe=white,
  boxrule=0pt,
  sharp corners,
  left=0mm,
  right=0mm,
  top=1mm,
  bottom=1mm,
  before skip=1mm,
  after skip=1mm,
  title={\color{edu-blue}\bfseries #1\hspace{1.5mm}#2},
  coltitle=edu-blue,
  attach title to upper={\par\vspace{2mm}},
  fontupper=\normalsize
}

\newenvironment{eduPlainList}{%
  \begin{enumerate}[
    label=(\arabic*),
    leftmargin=7mm,
    itemsep=1.2mm,
    topsep=0pt,
    parsep=0pt
  ]
}{%
  \end{enumerate}
}

% ---------- 兼容旧 block 的轻量盒子 ----------
\newtcolorbox{eduSoftBox}{
  enhanced,
  breakable,
  colback=edu-soft-bg,
  colframe=edu-line,
  boxrule=0.35pt,
  arc=1mm,
  left=3mm,
  right=3mm,
  top=2mm,
  bottom=2mm,
  before skip=1mm,
  after skip=2mm,
  fontupper=\small
}

\newtcolorbox{eduMistakeBox}{
  enhanced,
  breakable,
  colback=edu-red-bg,
  colframe=edu-red!40,
  boxrule=0.35pt,
  arc=1mm,
  left=3mm,
  right=3mm,
  top=2.5mm,
  bottom=2.5mm,
  before skip=1mm,
  after skip=2mm,
  fontupper=\normalsize
}

\newtcolorbox{eduLegacySolution}{
  enhanced,
  breakable,
  colback=edu-soft-bg,
  colframe=edu-line,
  boxrule=0.35pt,
  arc=1mm,
  left=3mm,
  right=3mm,
  top=2mm,
  bottom=2mm,
  before skip=-2mm,
  after skip=4mm,
  fontupper=\small
}

\newtcolorbox{eduTeacherNote}{
  enhanced,
  breakable,
  colback=edu-soft-bg,
  colframe=edu-line,
  boxrule=0pt,
  borderline west={1.5pt}{0pt}{edu-muted},
  left=2.5mm,
  right=2.5mm,
  top=2mm,
  bottom=2mm,
  before skip=1mm,
  after skip=2mm,
  fontupper=\footnotesize,
  colupper=edu-muted
}

% ---------- 答题区 ----------
\newcommand{\answerarea}[1][40mm]{%
  \vspace{3mm}%
  \begin{tcolorbox}[
    enhanced,
    breakable,
    colback=white,
    colframe=edu-line,
    boxrule=0.55pt,
    arc=0.8mm,
    height=#1,
    valign=top,
  ]
  \end{tcolorbox}%
}

\examsetup{
  question/show-answer = false,
  fillin/show-answer = false,
  solution/show-solution = hide,
}

\begin{document}


\begin{eduProblemBlock}[例 1]
（2018$\cdot$普陀区模拟）如图，是一个隧道的截面，如果路面 $AB$ 宽为 $8$ 米，净高 $CD$ 为 $8$ 米，那么这个隧道所在圆的半径 $OA$ 是\_\_\_\_米。
\end{eduProblemBlock}









\begin{eduAuxItem}{解题流程}
识别 $CD \perp AB$ 且 $CD$ 过圆心 $O$$\;\to\;$垂径定理得 $AD = 4$$\;\to\;$设半径 $r$，用 $r$ 表示 $OD$$\;\to\;$在 $\mathrm{Rt}\triangle AOD$ 中列方程，解出 $r$\end{eduAuxItem}




\begin{eduAuxItem}{关键想法}
净高 $CD$ 是从路面（弦 $AB$）到弧顶（圆弧最高点 $C$）的距离。
$CD$ 是一条垂直于 $AB$ 的线段，而且它经过圆心 $O$！

为什么？因为弧顶 $C$ 是圆上离弦 $AB$ 最远的点，过 $C$ 作 $AB$ 的垂线一定过圆心（垂径定理推论）。

所以 $C$、$O$、$D$ 三点共线，$CD$ 可以拆成两段：
\begin{itemize}
  \item $CO = r$（半径）
  \item $OD = CD - CO = 8 - r$
\end{itemize}
在 $\mathrm{Rt}\triangle AOD$ 中：$r^2 = 4^2 + (8-r)^2$。
\end{eduAuxItem}




\begin{eduSubQuestionItem}{第一步}
识别垂径关系并求半弦长。\end{eduSubQuestionItem}

\begin{eduDualExplain}
  \begin{eduExplainLeft}
    \eduColumnTitle{\eduIconBulb}{思考引导}
    \begin{eduNumberList}
      \item $CD$ 是从弧顶到路面的垂直距离——它和圆心有什么关系？      \item 弧顶 $C$ 是圆上到弦 $AB$ 最远的点。从 $C$ 向 $AB$ 作垂线，这条垂线过圆心吗？      \item 垂径定理推论：平分弧的直径垂直平分这条弧所对的弦。反过来想——从弧顶作弦的垂线，一定经过圆心。    \end{eduNumberList}
  \end{eduExplainLeft}
  \eduColumnDivider
  \begin{eduExplainRight}
    \eduColumnTitle{\eduIconBook}{规范讲解}
    \begin{eduNumberList}
      \item 设隧道截面为 $\odot O$，路面为弦 $AB$，弧顶为 $C$，$CD \perp AB$ 于 $D$。      \item $\because$ $C$ 是弧顶，$CD$ 是弧到弦的最大距离，$\therefore$ $CD$ 过圆心 $O$。      \item $\therefore$ $C$、$O$、$D$ 三点共线，$CD \perp AB$。      \item 由垂径定理：$D$ 是 $AB$ 中点，$AD = \frac{1}{2}AB = 4$。    \end{eduNumberList}

    \vspace{1mm}
    \eduSubTitle{关键}
    \begin{eduBulletList}
      \item 不要把 $CD=8$ 当成弦心距！弦心距是 $OD$，不是 $CD$。    \end{eduBulletList}
  \end{eduExplainRight}
\end{eduDualExplain}




\begin{eduSubQuestionItem}{第二步}
列方程求半径 $r$。\end{eduSubQuestionItem}

\begin{eduDualExplain}
  \begin{eduExplainLeft}
    \eduColumnTitle{\eduIconBulb}{思考引导}
    \begin{eduNumberList}
      \item $C$、$O$、$D$ 共线，$CD$ 从上到下依次是 $C \to O \to D$。那么 $OD$ 等于什么？      \item $CD = CO + OD = r + OD$，所以 $OD = 8 - r$。      \item 现在 $\mathrm{Rt}\triangle AOD$ 三边都知道了吗？哪条边还用 $r$ 表示？    \end{eduNumberList}
  \end{eduExplainLeft}
  \eduColumnDivider
  \begin{eduExplainRight}
    \eduColumnTitle{\eduIconBook}{规范讲解}
    \begin{eduNumberList}
      \item $\because$ $C$ 在 $O$ 上方，$D$ 在 $O$ 下方，$CD = CO + OD$。      \item $OD = CD - CO = 8 - r$。      \item 在 $\mathrm{Rt}\triangle AOD$ 中，由勾股定理：      \item $OA^2 = AD^2 + OD^2$      \item $r^2 = 4^2 + (8-r)^2 = 16 + 64 - 16r + r^2$      \item $16r = 80$，解得 $r = 5$。    \end{eduNumberList}

    \vspace{1mm}
    \eduSubTitle{注意}
    \begin{eduBulletList}
      \item $r^2$ 项消掉了！这题本质上是一元一次方程，不是二次方程。    \end{eduBulletList}
  \end{eduExplainRight}
\end{eduDualExplain}




\begin{eduBottomGrid}
  \begin{eduSummaryLeft}
    \begin{eduSummaryBlock}{\eduIconCheck}{答案}
    \begin{eduPlainList}
      \item $r = OA = 5$ 米    \end{eduPlainList}
    \end{eduSummaryBlock}
  \end{eduSummaryLeft}
  \eduSummaryDivider
  \begin{eduSummaryRight}
    \begin{eduSummaryBlock}{\eduIconPin}{方法提醒}
    \begin{eduBulletList}
      \item 「弦长 + 拱高求半径」的标准步骤：拱高拆为 $CO+OD$，设 $r$ 列方程。      \item $r^2$ 一定会消去，本质是一元一次方程。      \item 通用公式：$r = \frac{a^2 + 4h^2}{8h}$（$a$=弦长，$h$=拱高）。    \end{eduBulletList}
    \end{eduSummaryBlock}
  \end{eduSummaryRight}
\end{eduBottomGrid}




\begin{eduMistakeBox}
\eduSubTitle{把净高 $CD$ 当成弦心距 $OD$}弦心距是 $OD$（圆心到弦的距离），不是 $CD$（弦到弧顶的距离）。
$OD = CD - r$，必须先拆分。
\end{eduMistakeBox}




\begin{eduAuxItem}{思考}
如果把路面宽度改为 $6$ 米，净高改为 $2$ 米，半径是多少？试试用通用公式。\end{eduAuxItem}




\begin{eduAuxItem}{思考}
反过来：已知半径 $r=5$，弦长 $AB=8$，求净高（拱高）$CD$。\end{eduAuxItem}



\clearpage

\begin{eduAuxItem}{关键想法}
学生：陆子辰，八年级，B band。
本题定位：填空题，1 分钟。
预期能力：陆子辰能识别垂径定理，会设方程，应能在 1 分钟内完成。
讲解重点：强化「拱高 $\neq$ 弦心距」的区分，以及 $CD = CO + OD$ 的拆分。
如果陆子辰卡住：问他「$OD$ 用 $r$ 怎么表示？」
\end{eduAuxItem}




\begin{eduProblemBlock}[例 2]
（2021$\cdot$浦东新区模拟）如图，已知 $AB$ 是圆 $O$ 的直径，弦 $CD$ 交 $AB$ 于点 $E$，$\angle CEA = 30^\circ$，$OE = 4$，$DE = 5\sqrt{3}$，求弦 $CD$ 及圆 $O$ 的半径长。

\end{eduProblemBlock}









\begin{eduAuxItem}{解题流程}
过 $O$ 作 $OM \perp CD$ 于 $M$（核心辅助线）$\;\to\;$$30^\circ$ 特殊角求 $OM$ 和 $EM$$\;\to\;$分段求 $DM = DE - EM$$\;\to\;$垂径定理得 $CD$，勾股定理得半径\end{eduAuxItem}




\begin{eduAuxItem}{关键想法}
弦 $CD$ 不垂直于直径 $AB$，不能直接用垂径定理。
解决办法：自己制造垂直条件——过圆心 $O$ 作 $OM \perp CD$。

这样就有了两个直角三角形：
\begin{itemize}
  \item $\mathrm{Rt}\triangle OEM$：$\angle OEM = \angle CEA = 30^\circ$，$OE = 4$，可以求 $OM$ 和 $EM$
  \item $\mathrm{Rt}\triangle DOM$：$OM$ 已知，$DM$ 可算，用勾股定理求 $OD$（半径）
\end{itemize}
\end{eduAuxItem}




\begin{eduSubQuestionItem}{第一步}
作辅助线，利用 $30^\circ$ 特殊角求 $OM$ 和 $EM$。\end{eduSubQuestionItem}

\begin{eduDualExplain}
  \begin{eduExplainLeft}
    \eduColumnTitle{\eduIconBulb}{思考引导}
    \begin{eduNumberList}
      \item 见到弦 $CD$，想用垂径定理——但 $CD$ 不和直径垂直。怎么办？      \item 过 $O$ 作 $OM \perp CD$——这条辅助线你一定要形成条件反射！      \item 现在看 $\mathrm{Rt}\triangle OEM$：$\angle OEM$ 等于多少？为什么？      \item $30^\circ$ 的直角三角形，斜边 $OE=4$，你能快速说出 $OM$ 和 $EM$ 吗？    \end{eduNumberList}
  \end{eduExplainLeft}
  \eduColumnDivider
  \begin{eduExplainRight}
    \eduColumnTitle{\eduIconBook}{规范讲解}
    \begin{eduNumberList}
      \item 过圆心 $O$ 作 $OM \perp CD$ 于 $M$，联结 $OD$。      \item $\because \angle OEM = \angle CEA = 30^\circ$（对顶角相等）。      \item 在 $\mathrm{Rt}\triangle OEM$ 中，$OE = 4$：      \item $OM = \frac{1}{2}OE = \frac{1}{2} \times 4 = 2$（$30^\circ$ 所对的直角边 = 斜边的一半）。      \item $EM = OE \cdot \cos 30^\circ = 4 \times \frac{\sqrt{3}}{2} = 2\sqrt{3}$。    \end{eduNumberList}

    \vspace{1mm}
    \eduSubTitle{关键}
    \begin{eduBulletList}
      \item 30° 特殊角直角三角形三边比 $1 : \sqrt{3} : 2$ 要熟记。    \end{eduBulletList}
  \end{eduExplainRight}
\end{eduDualExplain}




\begin{eduSubQuestionItem}{第二步}
分段求 $DM$，再由垂径定理求 $CD$，勾股定理求半径。\end{eduSubQuestionItem}

\begin{eduDualExplain}
  \begin{eduExplainLeft}
    \eduColumnTitle{\eduIconBulb}{思考引导}
    \begin{eduNumberList}
      \item $M$ 是 $CD$ 的中点（垂径定理）。$E$ 在弦 $CD$ 上，$D$、$M$、$E$ 的顺序是怎样的？      \item $DE = 5\sqrt{3}$，$EM = 2\sqrt{3}$。$DM$ 等于 $DE + EM$ 还是 $DE - EM$？      \item 画一下：$M$ 是中点，$E$ 在 $M$ 的右边（靠近 $C$ 的方向），所以 $DM = DE - EM$。    \end{eduNumberList}
  \end{eduExplainLeft}
  \eduColumnDivider
  \begin{eduExplainRight}
    \eduColumnTitle{\eduIconBook}{规范讲解}
    \begin{eduNumberList}
      \item $\because OM \perp CD$ 且 $OM$ 过圆心，$\therefore M$ 是 $CD$ 中点。      \item $E$ 在线段 $DM$ 的延长线上（靠近 $C$），$\therefore DM = DE - EM$。      \item $DM = 5\sqrt{3} - 2\sqrt{3} = 3\sqrt{3}$。      \item 由垂径定理：$CD = 2DM = 2 \times 3\sqrt{3} = 6\sqrt{3}$。      \item 在 $\mathrm{Rt}\triangle DOM$ 中：      \item $OD = \sqrt{OM^2 + DM^2} = \sqrt{2^2 + (3\sqrt{3})^2} = \sqrt{4 + 27} = \sqrt{31}$。    \end{eduNumberList}

    \vspace{1mm}
    \eduSubTitle{注意}
    \begin{eduBulletList}
      \item 求半径一定要联结 $OD$（或 $OC$），在 $\mathrm{Rt}\triangle DOM$ 中用勾股定理。    \end{eduBulletList}
  \end{eduExplainRight}
\end{eduDualExplain}




\begin{eduBottomGrid}
  \begin{eduSummaryLeft}
    \begin{eduSummaryBlock}{\eduIconCheck}{答案}
    \begin{eduPlainList}
      \item 弦 $CD = 6\sqrt{3}$      \item 半径 $r = \sqrt{31}$    \end{eduPlainList}
    \end{eduSummaryBlock}
  \end{eduSummaryLeft}
  \eduSummaryDivider
  \begin{eduSummaryRight}
    \begin{eduSummaryBlock}{\eduIconPin}{方法提醒}
    \begin{eduBulletList}
      \item 见到弦 + 非垂直条件，第一反应：作弦心距（过圆心作弦的垂线段）。      \item 30° 特殊角直角三角形边比 $1:\sqrt{3}:2$ 要脱口而出。      \item 注意 $DM = DE - EM$ 的方向判断——画图确认 $M$ 和 $E$ 的位置。    \end{eduBulletList}
    \end{eduSummaryBlock}
  \end{eduSummaryRight}
\end{eduBottomGrid}




\begin{eduMistakeBox}
\eduSubTitle{$DM$ 方向搞反}$M$ 是 $CD$ 中点，$E$ 在弦上。如果搞不清 $D$、$M$、$E$ 的顺序，会把 $DM$ 算成 $DE + EM = 7\sqrt{3}$（错）。
正确做法：$DM = DE - EM = 3\sqrt{3}$。画图确认！
\end{eduMistakeBox}




\begin{eduAuxItem}{思考}
如果 $\angle CEA = 45^\circ$，其他条件不变，你会怎么算？和 $30^\circ$ 有什么区别？\end{eduAuxItem}




\begin{eduAuxItem}{思考}
为什么作 $OM \perp CD$ 后，$\angle OEM$ 就等于 $\angle CEA$？用了什么几何知识？\end{eduAuxItem}



\clearpage

\begin{eduAuxItem}{关键想法}
学生：陆子辰，八年级，B band，会作辅助线。
本题定位：解答题，5 分钟。
预期能力：陆子辰已有弦心距辅助线基础，应在 3-5 分钟内完成。
讲解重点：强化辅助线条件反射 + $DM = DE - EM$ 方向判断。
如果陆子辰卡住：问他「$M$ 在 $D$ 和 $E$ 之间，还是 $E$ 在 $D$ 和 $M$ 之间？」
\end{eduAuxItem}




\begin{eduProblemBlock}[例 3]
已知 $\odot O$ 半径 $OB = 13$，$AB$、$CD$ 是弦，$AB \parallel CD$，$AB = 10$，$CD = 24$，则 $AB$ 与 $CD$ 的距离是\_\_\_\_。
\end{eduProblemBlock}









\begin{eduAuxItem}{解题流程}
分别对 $AB$、$CD$ 作弦心距 $OM$、$ON$$\;\to\;$利用 $AB \parallel CD$ 证明 $OM$ 与 $ON$ 共线$\;\to\;$勾股定理分别求 $OM$ 和 $ON$$\;\to\;$分同侧（$OM - ON$）和异侧（$OM + ON$）求距离\end{eduAuxItem}




\begin{eduAuxItem}{关键想法}
两条平行弦 $AB \parallel CD$，从圆心 $O$ 分别向它们作垂线段 $OM$ 和 $ON$。
因为两条垂线都垂直于平行弦，所以 $OM \parallel ON$。
又因为它们都过点 $O$，所以 $M$、$O$、$N$ 三点共线！

两弦距离就等于弦心距的差或和：
\begin{itemize}
  \item 同侧：距离 $= OM - ON$
  \item 异侧：距离 $= OM + ON$
\end{itemize}
而 $OM$、$ON$ 各自独立用垂径定理 $+$ 勾股定理求。
\end{eduAuxItem}




\begin{eduSubQuestionItem}{第一步}
分别求两弦的弦心距 $OM$ 和 $ON$。\end{eduSubQuestionItem}

\begin{eduDualExplain}
  \begin{eduExplainLeft}
    \eduColumnTitle{\eduIconBulb}{思考引导}
    \begin{eduNumberList}
      \item 圆中有弦，过圆心作弦的垂线段——每条弦各作一条。      \item $AB \parallel CD$，$OM \perp AB$，$ON \perp CD$——两条垂线有什么关系？      \item 两条直线都垂直于平行线，它们互相平行。又都过点 $O$，所以三点共线！      \item 分别算 $OM$ 和 $ON$：半弦长分别是多少？    \end{eduNumberList}
  \end{eduExplainLeft}
  \eduColumnDivider
  \begin{eduExplainRight}
    \eduColumnTitle{\eduIconBook}{规范讲解}
    \begin{eduNumberList}
      \item 过 $O$ 作 $OM \perp AB$ 于 $M$，$ON \perp CD$ 于 $N$，联结 $OA$、$OC$。      \item $\because AB \parallel CD$，$OM \perp AB$，$ON \perp CD$，      \item $\therefore$ 点 $M$、$O$、$N$ 三点共线。      \item 由垂径定理：$AM = \frac{1}{2}AB = 5$，$CN = \frac{1}{2}CD = 12$。      \item 在 $\mathrm{Rt}\triangle OAM$ 中：$OM = \sqrt{OA^2 - AM^2} = \sqrt{13^2 - 5^2} = \sqrt{144} = 12$。      \item 在 $\mathrm{Rt}\triangle OCN$ 中：$ON = \sqrt{OC^2 - CN^2} = \sqrt{13^2 - 12^2} = \sqrt{25} = 5$。    \end{eduNumberList}

    \vspace{1mm}
    \eduSubTitle{关键}
    \begin{eduBulletList}
      \item 短弦 $AB=10$ 的弦心距反而更大（$OM=12$），长弦 $CD=24$ 的弦心距更小（$ON=5$）——弦越长，离圆心越近。    \end{eduBulletList}
  \end{eduExplainRight}
\end{eduDualExplain}




\begin{eduSubQuestionItem}{第二步}
分类讨论求两弦距离。\end{eduSubQuestionItem}

\begin{eduDualExplain}
  \begin{eduExplainLeft}
    \eduColumnTitle{\eduIconBulb}{思考引导}
    \begin{eduNumberList}
      \item 题目说两弦在圆心哪一侧了吗？没有！所以必须讨论两种情况。      \item 情形一：两弦在圆心同侧，距离怎么算？      \item 情形二：两弦在圆心异侧，距离怎么算？      \item 填空题答案要写几个？两个都要写！    \end{eduNumberList}
  \end{eduExplainLeft}
  \eduColumnDivider
  \begin{eduExplainRight}
    \eduColumnTitle{\eduIconBook}{规范讲解}
    \begin{eduNumberList}
      \item $M$、$O$、$N$ 共线，分两种情况：      \item 情形一（同侧）：$AB$ 和 $CD$ 在圆心同侧，距离 $= OM - ON = 12 - 5 = 7$。      \item 情形二（异侧）：$AB$ 和 $CD$ 在圆心异侧，距离 $= OM + ON = 12 + 5 = 17$。      \item $\therefore$ $AB$ 与 $CD$ 的距离为 $7$ 或 $17$。    \end{eduNumberList}

    \vspace{1mm}
    \eduSubTitle{注意}
    \begin{eduBulletList}
      \item 分类讨论的两种情况都要写！只写一个扣一半分。这是陆子辰需要特别注意的。    \end{eduBulletList}
  \end{eduExplainRight}
\end{eduDualExplain}




\begin{eduBottomGrid}
  \begin{eduSummaryLeft}
    \begin{eduSummaryBlock}{\eduIconCheck}{答案}
    \begin{eduPlainList}
      \item 距离为 $7$ 或 $17$    \end{eduPlainList}
    \end{eduSummaryBlock}
  \end{eduSummaryLeft}
  \eduSummaryDivider
  \begin{eduSummaryRight}
    \begin{eduSummaryBlock}{\eduIconPin}{方法提醒}
    \begin{eduBulletList}
      \item 两平行弦求距离 = 分别求弦心距 + 分类讨论同侧/异侧。      \item 分类讨论的触发条件：题目没说两弦在圆心哪一侧。      \item 规律：弦越长，弦心距越小；弦越短，弦心距越大。    \end{eduBulletList}
    \end{eduSummaryBlock}
  \end{eduSummaryRight}
\end{eduBottomGrid}




\begin{eduMistakeBox}
\eduSubTitle{只算一种情况漏解}题目没有明确两弦在圆心的哪一侧，就必须分同侧和异侧两种情况讨论。
填空题写两个答案，只写一个扣一半分。
\end{eduMistakeBox}




\begin{eduAuxItem}{思考}
如果 $AB = CD$（两弦等长），同侧距离是多少？这个结果说明什么？\end{eduAuxItem}




\begin{eduAuxItem}{思考}
如果半径改为 $r$，弦长分别为 $a$ 和 $b$，推导一般公式 $d = \sqrt{r^2 - (a/2)^2} \pm \sqrt{r^2 - (b/2)^2}$。\end{eduAuxItem}



\clearpage

\begin{eduAuxItem}{关键想法}
学生：陆子辰，八年级，B band。已知弱点：分类讨论意识弱。
本题定位：解答题，5 分钟。
教学重点：这是针对陆子辰弱点的专项训练题。
预期行为：陆子辰很可能只算出 7 或只算出 17，需要在他写完第一个答案后追问"还有没有别的可能？"
如果陆子辰只写了一个答案：不直接纠正，问"两弦能不能在圆心的另一侧？"
\end{eduAuxItem}




\end{document}
